\documentclass[11pt]{article}
\usepackage{amsmath,amssymb,amsthm,amsfonts}
\usepackage{hyperref}
\usepackage{geometry}
\geometry{margin=1in}

\title{Prime-Encoded Eigen Solvers as Categorical Prime Flows:\\
A Multiplicity-Theoretic and Tensor-Quantum Framework}
\author{Citizen Gardens \\ The Prime Materia Commons}
\date{Mayy 5, 2026}

\begin{document}
\maketitle

\begin{abstract}
This report presents a unified, multiplicity-theoretic, and categorical
framework for prime-encoded eigenvalue solvers, with particular focus
on a prime-weighted Lanczos method and its tensor/quantum extensions.
Building on a base preprint introducing prime-encoded eigenvalue
decomposition (PEED), we formalize:
(i) a category of prime-labelled Krylov modules for the Lanczos
recurrence,
(ii) an endofunctor capturing the prime-weighted Lanczos flow as a
``prime move'' on this category,
(iii) spectral and dynamical invariants as functors and natural
transformations, and
(iv) an analogous categorical treatment of the tensor/quantum layer,
including quantum phase estimation and prime-labelled tensor-network
representations of eigenstates.
We include an executive summary, mathematical definitions and
propositions, and a concrete code sketch illustrating how a
prime-weighted Lanczos prototype can be implemented in practice.
The goal is to provide a prior-art defensive publication around this
family of prime-indexed, recursively stable eigen-solvers and their
categorical formalization.
\end{abstract}
\newpage
\tableofcontents
\newpage
\section{Executive Summary}

\subsection{Conceptual Overview}

The underlying preprint introduces a prime-encoded eigenvalue
decomposition (PEED) for an $n\times n$ Hermitian matrix $A$:
\[
P(A) = \sum_{p\in \mathbb{P}_A} \alpha_p\, p\, A_p,
\]
where $\mathbb{P}_A$ is a finite set of primes, $\alpha_p\in\mathbb{R}$,
and $A_p$ are linear operators encoding ``interaction components''
associated with each prime.\cite{PrimeEigenSolvers} Eigenvalues are
also prime-encoded as
\[
\lambda_i = \prod_{j} p_j^{e_{ij}}
\]
with integer exponents $e_{ij}$ representing spectral multiplicity
structure.\cite{PrimeEigenSolvers}

The preprint further develops:
\begin{itemize}
  \item a prime-based QR iteration,
  \item a prime-weighted Lanczos method that inserts prime weights into
    the recurrence coefficients and the induced tridiagonal matrix,
  \item a recursive feedback update on eigenvalue estimates,
  \item quantum-inspired extensions including quantum phase estimation,
    and
  \item a tensor representation of eigenstates
    $\Psi(t)=\sum_i \lambda_i |p_i\rangle\otimes|e_i\rangle$.\cite{PrimeEigenSolvers}
\end{itemize}

In this report, we:
\begin{enumerate}
  \item Recast the prime-weighted Lanczos procedure as a \emph{prime
  flow} in a category of prime-labelled Krylov modules.
  \item Define an endofunctor that advances the Krylov subspace and
  records the prime weights, together with a natural transformation
  representing the stepwise inclusion.
  \item Identify invariants of this flow (trace, prime-weighted
  off-diagonal energy, coupling ratios, exponent signatures) as functors
  and natural quantities in the category.
  \item Develop an analogous category of prime-labelled tensor modules
  for the quantum/tensor layer, including a quantum evolution functor,
  a phase-estimation functor, and expectation-value functors for
  tensor-network observables.
\end{enumerate}

The resulting framework aligns with a multiplicity-theoretic viewpoint:
mathematical objects are recursively generated prime-labelled
configurations, and algorithms such as Lanczos become discrete flows on
prime-indexed multiplicity spaces with well-defined invariants.

\subsection{Novelty and Prior-Art Value}

The combination of the following elements appears to be novel and thus
worth documenting as prior art:
\begin{itemize}
  \item Explicit use of \emph{prime labels} at the level of:
    \begin{itemize}
      \item decomposition $P(A)=\sum_p \alpha_p p A_p$,
      \item eigenvalue encoding $\lambda_i=\prod_j p_j^{e_{ij}}$,
      \item recurrence coefficients $\beta_k p_k$ in Lanczos,
      \item tensor states $\Psi(t)=\sum_i \lambda_i |p_i\rangle\otimes|e_i\rangle$.
    \end{itemize}
  \item A \emph{categorical} formulation of the prime-weighted Lanczos
    method:
    \begin{itemize}
      \item objects: prime-labelled Krylov modules,
      \item morphisms: maps preserving prime identity structure,
      \item endofunctor: prime-weighted Lanczos step,
      \item natural transformation: canonical inclusion (prime move).
    \end{itemize}
  \item Spectral and dynamical invariants (trace, prime-weighted energy,
    coupling ratios, exponent signatures) expressed as functors or
    natural quantities within this category.
  \item A parallel categorical treatment for the tensor/quantum layer:
    prime-labelled tensor modules, quantum evolution functor, and
    phase-estimation functor.
\end{itemize}

Together, these give a coherent, mathematically explicit backbone for
prime-encoded eigenvalue solvers that is suitable for defensive
publication.

\section{Mathematical Overview}

\subsection{Prime-Encoding of Matrices and Eigenvalues}

We recall the basic setup from the preprint.\cite{PrimeEigenSolvers}
Let $A$ be an $n \times n$ Hermitian matrix; consider the eigenvalue
problem
\[
A v = \lambda v.
\]
The prime-encoding of $A$ is given by
\begin{equation}
P(A) = \sum_{p\in\mathbb{P}_A} \alpha_p\, p\, A_p,
\label{eq:prime-decomp}
\end{equation}
where $\mathbb{P}_A$ is a finite set of distinct primes, $\alpha_p \in
\mathbb{R}$, and $A_p$ are linear operators on the same space as
$A$.\cite{PrimeEigenSolvers}

Eigenvalues are encoded as prime products:
\begin{equation}
\lambda_i = \prod_{j} p_j^{e_{ij}},
\label{eq:lambda-prime-encoding}
\end{equation}
where $e_{ij} \in \mathbb{Z}$ are exponents encoding spectral
characteristics.\cite{PrimeEigenSolvers}

\subsection{Prime-Weighted Lanczos Method}

The prime-weighted Lanczos method constructs the Krylov subspace
\begin{equation}
K_m(A,v) = \mathrm{span}\{v,Av,A^2 v,\dots, A^{m-1}v\}
\label{eq:Krylov}
\end{equation}
and iteratively computes a tridiagonal matrix $T_m$.\cite{PrimeEigenSolvers}

Initialization (Eq.~(11)):
\begin{equation}
w_1 = A v_1,\quad \alpha_1 = v_1^T w_1.
\end{equation}

Iterative prime-weighted recurrence (Eq.~(12)-(13)):
\begin{align}
w_{k+1} &= A v_k - \alpha_k v_k - \beta_k v_{k-1}, \label{eq:lanczos-rec}\\
\beta_k &= \|w_k\| p_k, \label{eq:beta-prime}
\end{align}
with $\alpha_k = v_k^T A v_k$ (Eq.~(14)).\cite{PrimeEigenSolvers}

The resulting prime-encoded tridiagonal matrix (Eq.~(15)) is
\begin{equation}
T_m =
\begin{bmatrix}
\alpha_1 & \beta_1 p_1 & 0 & \dots & 0 \\
\beta_1 p_1 & \alpha_2 & \beta_2 p_2 & \dots & 0 \\
0 & \beta_2 p_2 & \alpha_3 & \dots & 0 \\
\vdots & \vdots & \vdots & \ddots & \beta_{m-1}p_{m-1} \\
0 & 0 & 0 & \beta_{m-1}p_{m-1} & \alpha_m
\end{bmatrix}.
\label{eq:Tm}
\end{equation}

\subsection{Recursive Feedback and Quantum/Tensor Layer}

Eigenvalues are further refined via a recursive feedback loop (Eq.~(8)):
\begin{equation}
\lambda_{t+1} = \lambda_t + \alpha_t \sum_i p_i e^{-\beta_i t},
\label{eq:feedback}
\end{equation}
with learning rate $\alpha_t$ and scaling parameters
$\beta_i$.\cite{PrimeEigenSolvers}

Quantum phase estimation is modeled via a unitary $U$ with
\begin{equation}
U|v\rangle = e^{i\lambda} |v\rangle,
\label{eq:QPE-unitary}
\end{equation}
as in Eq.~(9).\cite{PrimeEigenSolvers}

A tensor representation of eigenstates is given by
\begin{equation}
\Psi(t) = \sum_i \lambda_i |p_i\rangle \otimes |e_i\rangle,
\label{eq:tensor-state}
\end{equation}
as in Eq.~(10), realized in a tensor product Hilbert space
$\mathcal{I}\otimes\mathcal{E}$ where $\{|p_i\rangle\}$ is a
prime-labelled basis and $\{|e_i\rangle\}$ are eigenvectors.\cite{PrimeEigenSolvers}

\section{Categorical Framework for Prime-Weighted Lanczos}

\subsection{Category of Prime-Labelled Krylov Modules}

\subsubsection*{Base data}

Let $A$ be an $n\times n$ Hermitian matrix with prime encoding
$P(A)$ as in (\ref{eq:prime-decomp}).\cite{PrimeEigenSolvers} Fix a
starting vector $v_1 \in \mathcal{H}$ and define Krylov subspaces
$\mathcal{C}_m := K_m(A,v_1)$ as in (\ref{eq:Krylov}).\cite{PrimeEigenSolvers}

Let
\[
\mathcal{L} := \bigoplus_{p\in\mathbb{P}_A} \mathbb{R}\, e_p
\]
be the prime identity module, a free $\mathbb{R}$-module with basis
indexed by primes.\cite{PrimeEigenSolvers,CategoryModules}

\begin{definition}[Prime-labelled module]
An object of $\mathbf{PrimeMod}_A$ is a triple
$M = (\mathcal{C},\mathcal{L},\tau)$ where:
\begin{itemize}
  \item $\mathcal{C}\subseteq\mathcal{H}$ is a finite-dimensional
    subspace (typically $\mathcal{C}_m$),
  \item $\mathcal{L}$ is the prime identity module above,
  \item $\tau : \mathcal{L}\otimes\mathcal{C} \to \mathcal{H}$ is
    linear and satisfies
  \[
  \tau(e_p\otimes v) = A_p v \quad \forall p\in\mathbb{P}_A,\; v\in\mathcal{C}.
  \]
\end{itemize}
\end{definition}

\begin{definition}[Morphisms]
A morphism $\phi : M_1=(\mathcal{C}_1,\mathcal{L}_1,\tau_1)
\to M_2=(\mathcal{C}_2,\mathcal{L}_2,\tau_2)$ in
$\mathbf{PrimeMod}_A$ is a pair
$(\phi_{\mathcal{C}},\phi_{\mathcal{L}})$ where:
\begin{itemize}
  \item $\phi_{\mathcal{C}} : \mathcal{C}_1 \to \mathcal{C}_2$ is linear,
  \item $\phi_{\mathcal{L}} : \mathcal{L}_1 \to \mathcal{L}_2$ is linear
    and satisfies $\phi_{\mathcal{L}}(e_p) = c_p e_p$ for some
    $c_p\neq 0$,
  \item the diagram
  \[
  \begin{array}{ccc}
  \mathcal{L}_1\otimes\mathcal{C}_1 & \xrightarrow{\tau_1} & \mathcal{H} \\
  \downarrow \scriptstyle{\phi_{\mathcal{L}}\otimes\phi_{\mathcal{C}}} & &
  \downarrow \scriptstyle{\mathrm{id}_{\mathcal{H}}} \\
  \mathcal{L}_2\otimes\mathcal{C}_2 & \xrightarrow{\tau_2} & \mathcal{H}
  \end{array}
  \]
  commutes.
\end{itemize}
\end{definition}

It is straightforward to verify that $\mathbf{PrimeMod}_A$ is indeed a
category (identity and composition respect the commutativity
condition).\cite{CategoryModules}

\subsection{Prime-Weighted Lanczos Functor}

Fix a sequence $\vec{p}=(p_1,p_2,\dots)$ of primes from
$\mathbb{P}_A$, as used in the recurrence
(\ref{eq:lanczos-rec})--(\ref{eq:beta-prime}).\cite{PrimeEigenSolvers}

\begin{definition}[Prime-weighted Lanczos functor]
Define $\mathcal{F}_{\vec{p}} : \mathbf{PrimeMod}_A \to
\mathbf{PrimeMod}_A$ as follows.
\begin{itemize}
  \item On objects $M_m =
    (\mathcal{C}_m,\mathcal{L},\tau_m)$: let $M_{m+1} =
    (\mathcal{C}_{m+1},\mathcal{L},\tau_{m+1})$ where
    $\mathcal{C}_{m+1}=K_{m+1}(A,v_1)$ and $\tau_{m+1}$ is the unique
    extension of $\tau_m$ satisfying
    \[
    \tau_{m+1}(e_p\otimes v_{m+1}) = A_p v_{m+1}
    \]
    for $v_{m+1}$ produced by a single prime-weighted Lanczos step
    using $p_{m+1}$.\cite{PrimeEigenSolvers}
  \item On morphisms
    $\phi=(\phi_{\mathcal{C}},\phi_{\mathcal{L}}) :
    M_m \to N_m$, let
    $\mathcal{F}_{\vec{p}}(\phi)=(\phi'_{\mathcal{C}},\phi_{\mathcal{L}})$
    where $\phi'_{\mathcal{C}}$ extends $\phi_{\mathcal{C}}$ and sends
    $v_{m+1}$ to the vector obtained via the same numerical recurrence
    in $N_m$, with identical scalars
    $\alpha_k,\beta_k$ and prime $p_{m+1}$.\cite{PrimeEigenSolvers,LanczosSLEPc}
\end{itemize}
\end{definition}

\begin{proposition}[Lanczos step as natural transformation]
Let $\iota_{M_m} : M_m \to M_{m+1}$ be the canonical inclusion
\[
(\iota_{M_m})_{\mathcal{C}} : \mathcal{C}_m\hookrightarrow\mathcal{C}_{m+1},\quad
(\iota_{M_m})_{\mathcal{L}} = \mathrm{id}_{\mathcal{L}}.
\]
Then the family $\{\iota_{M_m}\}$ is a natural transformation
\[
\iota : \mathrm{Id}_{\mathbf{PrimeMod}_A} \Rightarrow \mathcal{F}_{\vec{p}}.
\]
In the basis induced by iterating $\iota$ and $\mathcal{F}_{\vec{p}}$,
the matrix of $A$ on $\mathcal{C}_m$ is precisely the tridiagonal
matrix $T_m$ with diagonal entries $\alpha_k$ and off-diagonals
$\beta_k p_k$ as in Eq.~(\ref{eq:Tm}).\cite{PrimeEigenSolvers,LanczosSLEPc}
\end{proposition}

\subsection{Invariants as Functors and Natural Quantities}

\subsubsection{Trace}

\begin{definition}[Trace functor]
Define $\mathrm{Tr} : \mathbf{PrimeMod}_A \to \mathbb{R}$ by
\[
\mathrm{Tr}(M_m) = \sum_{k=1}^m \alpha_k,
\]
where $\alpha_k$ are the Lanczos diagonal coefficients for
$\mathcal{C}_m$.\cite{PrimeEigenSolvers,LanczosSLEPc}
\end{definition}

For isometric morphisms that intertwine $A$, $\mathrm{Tr}$ is
invariant; under $\iota$,
\[
\mathrm{Tr}(\mathcal{F}_{\vec{p}}(M_m)) = \mathrm{Tr}(M_m) + \alpha_{m+1}.
\]

\subsubsection{Prime-Weighted Off-Diagonal Energy}

\begin{definition}[Prime-weighted off-diagonal energy]
Define
\[
E(M_m) = \sum_{k=1}^{m-1} (\beta_k p_k)^2,
\]
the squared Frobenius norm of the off-diagonal part of
$T_m$.\cite{PrimeEigenSolvers,LanczosSLEPc}
\end{definition}

$E$ is invariant under morphisms preserving the products
$\beta_k p_k$; under $\iota$,
\[
E(\mathcal{F}_{\vec{p}}(M_m)) = E(M_m) + (\beta_m p_m)^2.
\]

\subsubsection{Coupling Ratios and Exponent Signature}

\begin{definition}[Scale-free coupling ratios]
For $m\ge2$, define
\[
r_k = \frac{\beta_k p_k}{\beta_{k-1} p_{k-1}},\quad k=2,\dots,m.
\]
\end{definition}

These ratios are invariant under \emph{prime-similarities}:
morphisms that rescale all $e_p$ by the same constant and preserve all
$\|w_k\|$.

\begin{definition}[Prime-exponent signature]
Define
\[
s_m = \sum_{k=1}^{m-1} \log_{p_k} |\beta_k p_k|.
\]
\end{definition}

Under morphisms that preserve $\beta_k p_k$ up to prime powers, $s_m$
transforms linearly; once the Lanczos process has converged on a
spectral block, $s_m$ stabilizes up to bounded variation.

\subsection{Recursive Feedback as a Derived Functor}

Let $\mathbf{Dyn}(\mathbb{R})$ denote the category of discrete-time
dynamical systems on $\mathbb{R}$.

\begin{definition}[Feedback functor]
Given $M_m$ with primes and scalars $(\alpha_\tau,\beta_i)$, define
\[
\mathcal{R}(M_m)(t) = \lambda_0 +
\sum_{\tau=0}^{t-1}
\alpha_\tau \sum_{i} p_i e^{-\beta_i \tau}
\]
with initialization $\lambda_0$. This yields a map
$\mathcal{R} : \mathbf{PrimeMod}_A \to
\mathbf{Dyn}(\mathbb{R})$.
\end{definition}

Morphisms preserving the scalar data induce identical dynamical systems
under $\mathcal{R}$, making the feedback update functorial.

\section{Categorical Framework for Tensor/Quantum Layer}

\subsection{Category of Prime-Labelled Tensor Modules}

\subsubsection*{Base data}

Let $A$ have prime encoding $P(A)$ as in (\ref{eq:prime-decomp}) and
eigen-decomposition $A |e_i\rangle = \lambda_i |e_i\rangle$ with
prime encodings $\lambda_i$ as in (\ref{eq:lambda-prime-encoding}).\cite{PrimeEigenSolvers}
Let $\mathcal{I}$ be a Hilbert space with orthonormal basis
$\{|p\rangle : p\in\mathbb{P}_A\}$ and $\mathcal{E}$ the eigenspace
spanned by $\{|e_i\rangle\}$.\cite{PrimeEigenSolvers}

\begin{definition}[Prime-tensor module]
An object of $\mathbf{PrimeTen}_A$ is a triple
$N=(\mathcal{E},\mathcal{I},\Theta)$ where:
\begin{itemize}
  \item $\mathcal{E}\subseteq\mathcal{H}$ is a finite-dimensional
    eigenspace,
  \item $\mathcal{I}$ is the prime identity Hilbert space,
  \item $\Theta : \mathcal{I}\otimes\mathcal{E} \to \mathcal{H}$ is
    linear and satisfies
    \[
    \Theta(|p_i\rangle\otimes|e_i\rangle) = \lambda_i |e_i\rangle
    \]
    for each eigenpair $(\lambda_i,|e_i\rangle)$.
\end{itemize}
\end{definition}

The prime-labelled tensor state $\Psi(t)$ from
(\ref{eq:tensor-state}) lies naturally in $\mathcal{I}\otimes\mathcal{E}$.

\begin{definition}[Morphisms]
A morphism $\Phi : N_1=(\mathcal{E}_1,\mathcal{I}_1,\Theta_1)
\to N_2=(\mathcal{E}_2,\mathcal{I}_2,\Theta_2)$ is a pair
$(\Phi_{\mathcal{E}},\Phi_{\mathcal{I}})$ where:
\begin{itemize}
  \item $\Phi_{\mathcal{E}} : \mathcal{E}_1 \to \mathcal{E}_2$ is
    linear and intertwines $A$,
  \item $\Phi_{\mathcal{I}} : \mathcal{I}_1 \to \mathcal{I}_2$ is
    linear and satisfies $\Phi_{\mathcal{I}}(|p\rangle)
    = e^{i\theta_p}|p\rangle$,
  \item the diagram
  \[
  \begin{array}{ccc}
  \mathcal{I}_1\otimes\mathcal{E}_1 & \xrightarrow{\Theta_1} & \mathcal{H} \\
  \downarrow \scriptstyle{\Phi_{\mathcal{I}}\otimes\Phi_{\mathcal{E}}} & &
  \downarrow \scriptstyle{\mathrm{id}_{\mathcal{H}}} \\
  \mathcal{I}_2\otimes\mathcal{E}_2 & \xrightarrow{\Theta_2} & \mathcal{H}
  \end{array}
  \]
  commutes.
\end{itemize}
\end{definition}

\subsection{Quantum Evolution and Phase Estimation Functors}

Let $U=e^{iA}$ be the unitary such that
$U|e_i\rangle = e^{i\lambda_i}|e_i\rangle$ as in Eq.~(9).\cite{PrimeEigenSolvers}

\begin{definition}[Quantum evolution functor]
Define $\mathcal{Q} : \mathbf{PrimeTen}_A \to
\mathbf{PrimeTen}_A$ by:
\begin{itemize}
  \item on objects $N=(\mathcal{E},\mathcal{I},\Theta)$,
    \[
    \mathcal{Q}(N) = (\mathcal{E},\mathcal{I},\Theta_U)
    \]
    with
    \[
    \Theta_U(|p_i\rangle\otimes|e_i\rangle) =
    e^{i\lambda_i}\lambda_i |e_i\rangle,
    \]
  \item on morphisms, $\mathcal{Q}(\Phi) = (\Phi_{\mathcal{E}},\Phi_{\mathcal{I}})$,
    since $\Phi_{\mathcal{E}}$ intertwines $A$ and thus $U$.
\end{itemize}
\end{definition}

Let $\mathbf{QCirc}$ denote a suitable category of quantum circuits and
measurement outcomes.

\begin{definition}[Phase-estimation functor]
A phase-estimation functor $\mathcal{P} :
\mathbf{PrimeTen}_A \to \mathbf{QCirc}$ assigns to each
$N=(\mathcal{E},\mathcal{I},\Theta)$ a quantum circuit implementing
phase estimation for $U$ on states in $\mathcal{E}$, with measurement
outcomes labelled by the prime basis states $|p_i\rangle$.
\end{definition}

\begin{proposition}[Functoriality of phase estimation]
For any morphism $\Phi : N_1\to N_2$ in $\mathbf{PrimeTen}_A$ that
intertwines $U$, the corresponding circuits $\mathcal{P}(N_1)$ and
$\mathcal{P}(N_2)$ produce identical phase distributions up to
relabelling via $\Phi_{\mathcal{I}}$.
\end{proposition}

\subsection{Tensor-Network Observables and Invariants}

The tensor state $\Psi(t)$ admits tensor-network representations such as
matrix product states.\cite{PrimeEigenSolversTensor} Prime labels may
appear on physical or bond indices.

\begin{definition}[Expectation-value functor]
Let $\mathcal{O}$ be an observable on $\mathcal{I}\otimes\mathcal{E}$. Define
$\mathbb{E}_{\mathcal{O}} : \mathbf{PrimeTen}_A \to \mathbb{R}$ by
\[
\mathbb{E}_{\mathcal{O}}(N) =
\langle \Psi(t)|\mathcal{O}|\Psi(t)\rangle,
\]
where $\Psi(t)$ is determined by $\Theta$.
\end{definition}

When $\Phi$ preserves $\mathcal{O}$ (up to conjugation),
$\mathbb{E}_{\mathcal{O}}(N)$ is invariant, providing a tensor-network
analogue of the Lanczos invariants (trace, $E(M_m)$, $r_k$, $s_m$).

\section{Code Sketch: Prime-Weighted Lanczos Prototype}

This section illustrates how a research prototype of the prime-weighted
Lanczos method can be implemented in Python/NumPy. The code is
presented as an illustrative sketch, not a production-ready solver.

\subsection{Prime-Weighted Lanczos in Python}

\begin{verbatim}
import numpy as np

def prime_weighted_lanczos(A, v1, primes, m_max):
    """
    Prime-weighted Lanczos prototype.

    Parameters
    ----------
    A : (n, n) Hermitian numpy array
    v1 : (n,) initial vector, will be normalized
    primes : list or array of primes [p_1, p_2, ...]
    m_max : maximum Krylov depth

    Returns
    -------
    alphas : array of length m
    betas  : array of length m-1
    T      : (m, m) tridiagonal matrix with off-diagonals beta_k * p_k
    V      : (n, m) Lanczos basis vectors
    """

    n = A.shape[0]
    v1 = v1 / np.linalg.norm(v1)
    V = np.zeros((n, m_max), dtype=A.dtype)
    V[:, 0] = v1

    alphas = np.zeros(m_max, dtype=float)
    betas = np.zeros(m_max - 1, dtype=float)

    w = A @ V[:, 0]
    alphas[0] = np.vdot(V[:, 0], w).real

    for k in range(m_max - 1):
        if k > 0:
            w = A @ V[:, k] - alphas[k] * V[:, k] - betas[k-1] * V[:, k-1]
        else:
            w = A @ V[:, k] - alphas[k] * V[:, k]

        # prime-weighted beta
        p_k = primes[k]
        beta_k_norm = np.linalg.norm(w)
        betas[k] = beta_k_norm * p_k

        if beta_k_norm == 0:
            # early termination
            m = k + 1
            alphas = alphas[:m]
            betas = betas[:m-1]
            V = V[:, :m]
            break

        V[:, k+1] = w / beta_k_norm
        alphas[k+1] = np.vdot(V[:, k+1], A @ V[:, k+1]).real

    # construct T
    m = V.shape[1]
    T = np.diag(alphas[:m])
    for k in range(m - 1):
        p_k = primes[k]
        off = betas[k]  # already includes p_k
        T[k, k+1] = off
        T[k+1, k] = off

    return alphas[:m], betas[:m-1], T, V
\end{verbatim}

This prototype:
\begin{itemize}
  \item Implements the recurrence (\ref{eq:lanczos-rec})--(\ref{eq:beta-prime})
    with $\beta_k = \|w_k\| p_k$.
  \item Constructs the tridiagonal matrix $T_m$ of
    (\ref{eq:Tm}) with off-diagonals $\beta_k p_k$ embedded in
    \texttt{betas}.
  \item Produces the Lanczos basis $V$ for Krylov subspaces
    $\mathcal{C}_m$.
\end{itemize}

Extensions can:
\begin{itemize}
  \item track invariants $E(M_m)$, $r_k$, $s_m$ at each iteration,
  \item plug $T_m$ into a tensor-network representation (e.g., MPS/MPO),
  \item connect to a symbolic or categorical library that reflects
    the $\mathbf{PrimeMod}_A$ structure directly.
\end{itemize}

\section{Conclusions and Further Work}

This report consolidates the developments in prime-encoded eigenvalue
solvers into a unified multiplicity-theoretic and categorical
framework. The prime-weighted Lanczos method is cast as a prime flow in
$\mathbf{PrimeMod}_A$ with well-defined invariants, and the tensor/
quantum layer is described in $\mathbf{PrimeTen}_A$ with quantum
evolution and phase-estimation functors.

Future work could extend these ideas to:
\begin{itemize}
  \item Category-theoretic treatments of the prime-based QR iteration,
  \item Hybrid quantum-classical prime flows, where classical Lanczos
    and quantum phase estimation are combined as interconnected
    functors,
  \item Practical benchmarks comparing prime-encoded solvers with
    standard methods in high-dimensional settings.
\end{itemize}

\appendix
\section*{Mathematical Appendix}
\addcontentsline{toc}{section}{Mathematical Appendix}

This appendix collects explicit proofs, norm bounds, and auxiliary
lemmas underpinning the prime-weighted Lanczos framework and its
tensor/quantum extensions. We focus on:
\begin{itemize}
  \item well-posedness and stability of the prime-weighted Lanczos
    recurrence,
  \item operator-norm bounds for the induced tridiagonal $T_m$ and
    associated residuals,
  \item invariance properties of the prime-labelled categorical
    structures,
  \item compatibility of the tensor representation with quantum
    evolution and phase estimation.
\end{itemize}

Throughout, $A\in\mathbb{C}^{n\times n}$ is Hermitian and prime-encoded
as in Eq.~(2), and the Lanczos algorithm is as in Eqs.~(11)--(15) of
the main text.\cite{PrimeEigenSolvers,LanczosTour,UrschelLanczos}

\section{Prime-Weighted Lanczos: Well-Posedness and Norm Bounds}

\subsection{Well-Posedness of the Prime-Weighted Recurrence}

Recall the prime-weighted recurrence (Eqs.~(11)--(14)):
\begin{align}
w_1 &= A v_1, \qquad \alpha_1 = v_1^* w_1, \\
w_{k+1} &= A v_k - \alpha_k v_k - \beta_k v_{k-1}, \label{eq:app-lanczos-rec}\\
\beta_k &= \|w_k\|\, p_k, \label{eq:app-beta}\\
\alpha_k &= v_k^* A v_k,
\end{align}
with $v_0:=0$ and $p_k$ the $k$-th prime in a fixed sequence
$\vec{p}$.\cite{PrimeEigenSolvers}

\begin{lemma}[Nondegeneracy and early termination]
Assume $v_1\neq 0$ and $A$ Hermitian.
\begin{enumerate}
  \item If $w_k\neq 0$ for $k=1,\dots,m$, then the recurrence uniquely
    defines nonzero $\beta_k$ and unit vectors $v_{k+1}$ up to phase.
  \item If $w_{k_0}=0$ for some minimal $k_0$, the algorithm terminates
    with $\mathcal{C}_{k_0}=K_{k_0}(A,v_1)$ invariant under $A$ and
    exact on that subspace.
\end{enumerate}
\end{lemma}

\begin{proof}
(1) For a given $k$, if $w_k\neq 0$, then $\|w_k\|>0$ and hence
$\beta_k = \|w_k\| p_k \neq 0$ since $p_k>0$. Defining
$v_{k+1} = w_{k+1}/\|w_{k+1}\|$ produces a unit vector uniquely
determined up to a complex phase factor, as in the standard Lanczos
algorithm.\cite{LanczosTour,BerkeleyLanczos}

(2) If $w_{k_0}=0$, the recurrence gives $A v_{k_0} = \alpha_{k_0}
v_{k_0} + \beta_{k_0} v_{k_0-1}$; since $w_{k_0}=0$ implies
$\|w_{k_0}\|=0$, we have $\beta_{k_0}=0$ and thus
$A v_{k_0} = \alpha_{k_0} v_{k_0}$, so $v_{k_0}$ is an eigenvector.
Standard arguments for Lanczos show that
$\mathrm{span}\{v_1,\dots,v_{k_0}\}$ is $A$-invariant and that the
restriction of $A$ to $\mathcal{C}_{k_0}$ is represented exactly by
the tridiagonal $T_{k_0}$.\cite{LanczosTour,UrschelLanczos}
\end{proof}

\subsection{Operator Norm Bounds for the Tridiagonal Matrix}

Let $Q_m=[v_1,\dots,v_m]\in\mathbb{C}^{n\times m}$ be the orthonormal
basis of $\mathcal{C}_m$ generated by the recurrence. We have the
tridiagonal reduction
\begin{equation}
Q_m^* A Q_m = T_m
\end{equation}
with $T_m$ as in Eq.~(15).\cite{PrimeEigenSolvers,LanczosTour}

\begin{lemma}[Spectral inclusion]
For Hermitian $A$, we have
\[
\sigma(T_m)\subset \mathrm{conv}(\sigma(A)),
\]
where $\sigma(\cdot)$ denotes spectrum and $\mathrm{conv}$ the convex
hull.
\end{lemma}

\begin{proof}
This is standard for Lanczos: $T_m$ is the representation of $A$ on the
Krylov subspace with an orthonormal basis, hence its eigenvalues are
Ritz values for $A$ and lie in the convex hull of $\sigma(A)$ by the
variational characterization of eigenvalues and the Rayleigh quotient
bounds.\cite{LanczosAlgo,WikiLanczos}
\end{proof}

\begin{lemma}[Spectral norm bound]
Let $\|\cdot\|_2$ denote the spectral norm. Then
\[
\|T_m\|_2 \le \|A\|_2.
\]
\end{lemma}

\begin{proof}
Since $T_m = Q_m^* A Q_m$ with $Q_m$ having orthonormal columns,
\[
\|T_m\|_2 = \|Q_m^* A Q_m\|_2 \le \|Q_m^*\|_2 \|A\|_2 \|Q_m\|_2
= \|A\|_2,
\]
as $\|Q_m\|_2 = \|Q_m^*\|_2 = 1$.\cite{LanczosAlgo}
\end{proof}

The prime weights appear only in the off-diagonal entries
$\beta_k p_k$. A crude but explicit bound in terms of primes is:

\begin{proposition}[Prime-weighted Gershgorin bound]
Let $T_m$ be as in Eq.~(15). Then every eigenvalue $\theta\in\sigma(T_m)$
satisfies
\[
|\theta - \alpha_k|\le |\beta_{k-1} p_{k-1}| + |\beta_k p_k|
\]
for some $k$, with the convention $\beta_0 p_0 := 0$. Consequently,
\[
\max_{\theta\in\sigma(T_m)} |\theta|
\le \max_k \left( |\alpha_k|+|\beta_{k-1} p_{k-1}|+|\beta_k p_k| \right).
\]
\end{proposition}

\begin{proof}
This follows from Gershgorin's circle theorem applied to the
tridiagonal matrix $T_m$.\cite{UrschelLanczos} The $k$-th row has
diagonal $\alpha_k$ and off-diagonal entries at most
$|\beta_{k-1}p_{k-1}|$ and $|\beta_k p_k|$ in magnitude. The spectrum
lies in the union of those disks, yielding the bound.
\end{proof}

\subsection{Residual Norm Bounds and Ritz Error Estimates}

Let $\theta^{(m)}_j$ be an eigenvalue of $T_m$ (a Ritz value of $A$),
with associated Ritz vector $u^{(m)}_j = Q_m y_j$ where $y_j$ is a unit
eigenvector of $T_m$. The residual for this Ritz pair is
\[
r^{(m)}_j = A u^{(m)}_j - \theta^{(m)}_j u^{(m)}_j.
\]

\begin{lemma}[Residual norm formula]
For Hermitian $A$ and tridiagonal $T_m$, we have
\[
\|r^{(m)}_j\|_2 = |\beta_m p_m|\, |e_m^T y_j|,
\]
where $e_m$ is the $m$-th standard basis vector in $\mathbb{R}^m$.
\end{lemma}

\begin{proof}
In the classical Lanczos setting, the residual norm is
$\|r^{(m)}_j\|_2 = |\beta_m| |e_m^T y_j|$.\cite{LanczosAlgo,MITLanczos}
In the prime-weighted case, $T_m$ has off-diagonal entries $\beta_k p_k$,
and the coupling to the $(m+1)$-st basis vector in the recurrence is
scaled accordingly. Direct computation of
$A Q_m - Q_m T_m$ gives the residual term at the last basis vector
proportional to $\beta_m p_m v_{m+1}$; projecting onto $y_j$ yields the
stated expression.
\end{proof}

\begin{proposition}[A priori Ritz error bound]
Let $\lambda_1\ge\lambda_2\ge\dots\ge\lambda_n$ be the eigenvalues of
$A$, and let $\theta^{(m)}_1$ be the largest Ritz value of $T_m$.
Then
\[
\lambda_1 - \theta^{(m)}_1 \le \|r^{(m)}_1\|_2.
\]
In particular,
\[
\lambda_1 - \theta^{(m)}_1 \le |\beta_m p_m|\, |e_m^T y_1|.
\]
\end{proposition}

\begin{proof}
This is a standard bound for symmetric Lanczos: the Ritz value
$\theta^{(m)}_1$ is the maximum of the Rayleigh quotient on the Krylov
subspace, and the residual norm bounds the distance to a true
eigenpair.\cite{LanczosAlgo,UrschelLanczos} The second inequality comes
from the previous lemma.
\end{proof}

These bounds show how the prime factors $p_m$ enter the residual error
and thus affect convergence; they also motivate normalisation or
rescaling strategies if $p_m$ grows too quickly.

\section{Categorical Invariants and Natural Transformations}

\subsection{Trace and Energy as Functors}

We recall the trace functor
$\mathrm{Tr}(M_m)=\sum_{k=1}^m\alpha_k$ and the energy
$E(M_m)=\sum_{k=1}^{m-1}(\beta_k p_k)^2$.\cite{PrimeEigenSolvers}

\begin{proposition}[Functoriality of trace]
Let $\phi : M_m\to N_m$ be a morphism in $\mathbf{PrimeMod}_A$ such
that $\phi_{\mathcal{C}}$ is an isometric isomorphism onto its image
and intertwines $A$. Then
\[
\mathrm{Tr}(M_m) = \mathrm{Tr}(N_m).
\]
\end{proposition}

\begin{proof}
Since $\phi_{\mathcal{C}}$ is isometric and intertwines $A$, it induces
a unitary similarity between the restrictions $A|_{\mathcal{C}_m}$ and
$A|_{\phi_{\mathcal{C}}(\mathcal{C}_m)}$, hence their traces
coincide. The Lanczos tridiagonals $T_m$ and $T_m'$ are unitarily
similar and share the same diagonal sum.\cite{LanczosAlgo}
\end{proof}

\begin{proposition}[Prime-weighted energy monotonicity]
Under the inclusion $\iota_{M_m} : M_m\to M_{m+1}$, we have
\[
E(M_{m+1}) = E(M_m) + (\beta_m p_m)^2 \ge E(M_m).
\]
\end{proposition}

\begin{proof}
By definition,
\[
E(M_{m+1}) = \sum_{k=1}^{m}(\beta_k p_k)^2 =
\sum_{k=1}^{m-1}(\beta_k p_k)^2 +(\beta_m p_m)^2 = E(M_m)+(\beta_m p_m)^2.
\]
\end{proof}

\subsection{Scale-Free Ratios and Exponent Signature}

\begin{proposition}[Invariance under prime-similarity]
Let $\phi : M_m\to N_m$ be a morphism in $\mathbf{PrimeMod}_A$
such that:
\begin{itemize}
  \item $\phi_{\mathcal{L}}(e_p) = c e_p$ for the same $c\neq 0$ and
    all primes $p$,
  \item $\phi_{\mathcal{C}}$ preserves all $\|w_k\|$.
\end{itemize}
Then the scale-free ratios
$r_k=(\beta_k p_k)/(\beta_{k-1} p_{k-1})$ are invariant.
\end{proposition}

\begin{proof}
Under the stated assumptions, both $\beta_k$ and $p_k$ are scaled
in such a way that the products $\beta_k p_k$ are multiplied by a
common factor, which cancels in the ratio. Preservation of $\|w_k\|$
ensures that the norm-derived parts of $\beta_k$ are unchanged; the
global scaling factor $c$ affects all primes equally and cancels in
the ratio.\cite{PrimeEigenSolvers}
\end{proof}

\begin{proposition}[Exponent signature stability]
Assume the Lanczos process has converged on an invariant subspace
corresponding to a cluster of eigenvalues $\{\lambda_i\}$.
Then the quantity
\[
s_m = \sum_{k=1}^{m-1} \log_{p_k}|\beta_k p_k|
\]
stabilizes as $m$ increases (up to bounded fluctuations).
\end{proposition}

\begin{proof}[Sketch]
Once the Krylov subspace captures the invariant subspace, the recursion
coefficients $(\alpha_k,\beta_k)$ converge to asymptotic values that
depend only on the spectral block.\cite{LanczosAlgo,ConvergenceKrylov}
The additional prime factors $p_k$ are fixed by the chosen sequence,
hence the growth of $|\beta_k p_k|$ becomes regular, and their log
sum behaves like a linear function of $m$ plus bounded fluctuations.
The precise constant depends on the asymptotic values of $\beta_k$ and
the distribution of $p_k$.
\end{proof}

\section{Tensor/Quantum Layer: Compatibility and Bounds}

\subsection{Compatibility of $\Theta$ with $A$ and $U$}

In the tensor/quantum category $\mathbf{PrimeTen}_A$, objects
$N=(\mathcal{E},\mathcal{I},\Theta)$ satisfy
$\Theta(|p_i\rangle\otimes|e_i\rangle) = \lambda_i |e_i\rangle$.

\begin{lemma}[Intertwining with $A$]
For any $N$ in $\mathbf{PrimeTen}_A$, we have
\[
A\,\Theta(|p_i\rangle\otimes|e_i\rangle)
= \lambda_i\,\Theta(|p_i\rangle\otimes|e_i\rangle).
\]
\end{lemma}

\begin{proof}
By definition,
$\Theta(|p_i\rangle\otimes|e_i\rangle) = \lambda_i |e_i\rangle$ and
$A|e_i\rangle = \lambda_i |e_i\rangle$. Thus
\[
A\,\Theta(|p_i\rangle\otimes|e_i\rangle) =
A(\lambda_i |e_i\rangle)=\lambda_i^2 |e_i\rangle =
\lambda_i\,\Theta(|p_i\rangle\otimes|e_i\rangle),
\]
so $\Theta$ acts as a prime-labelled eigenvalue insertion operator.
\end{proof}

Let $U=e^{iA}$ be the unitary used in phase estimation.\cite{PrimeEigenSolvers}

\begin{lemma}[Intertwining with $U$]
We have
\[
U\,\Theta(|p_i\rangle\otimes|e_i\rangle) =
e^{i\lambda_i}\lambda_i |e_i\rangle =
\Theta_U(|p_i\rangle\otimes|e_i\rangle),
\]
where $\Theta_U$ defines the quantum evolution functor
$\mathcal{Q}$ in the main text.
\end{lemma}

\begin{proof}
Since $U|e_i\rangle = e^{i\lambda_i} |e_i\rangle$,
\[
U\,\Theta(|p_i\rangle\otimes|e_i\rangle) =
U(\lambda_i |e_i\rangle) = e^{i\lambda_i}\lambda_i |e_i\rangle,
\]
which is precisely the action of $\Theta_U$ on the same basis element.
\end{proof}

\subsection{Norm Bounds for Tensor States}

The tensor state
$\Psi(t) = \sum_i \lambda_i |p_i\rangle\otimes|e_i\rangle$ satisfies:

\begin{lemma}[Norm of prime-labelled tensor state]
Assume $\{|p_i\rangle\}$ and $\{|e_i\rangle\}$ are orthonormal bases.
Then
\[
\|\Psi(t)\|^2 = \sum_i |\lambda_i|^2.
\]
\end{lemma}

\begin{proof}
Orthogonality of the tensor product basis implies
\[
\|\Psi(t)\|^2 =
\left\langle \sum_i \lambda_i |p_i,e_i\rangle,
\sum_j \lambda_j |p_j,e_j\rangle \right\rangle
=
\sum_i |\lambda_i|^2,
\]
where $|p_i,e_i\rangle := |p_i\rangle\otimes|e_i\rangle$.
\end{proof}

\begin{proposition}[Bound in terms of $\|A\|_2$]
If $A$ is Hermitian with eigenvalues $\lambda_i$ and spectral norm
$\|A\|_2=\max_i |\lambda_i|$, then
\[
\|\Psi(t)\|^2 \le n \|A\|_2^2.
\]
\end{proposition}

\begin{proof}
We have $|\lambda_i|\le \|A\|_2$ for all $i$, so
\[
\|\Psi(t)\|^2 = \sum_{i=1}^n |\lambda_i|^2 \le n \|A\|_2^2.
\]
\end{proof}

\subsection{Expectation-Value Bounds for Tensor Observables}

Let $\mathcal{O}$ be a bounded operator on $\mathcal{I}\otimes\mathcal{E}$.

\begin{lemma}[Expectation-value bound]
We have
\[
|\langle \Psi(t)|\mathcal{O}|\Psi(t)\rangle|
\le \|\mathcal{O}\|_2\, \|\Psi(t)\|^2.
\]
\end{lemma}

\begin{proof}
By Cauchy--Schwarz,
\[
|\langle \Psi(t)|\mathcal{O}|\Psi(t)\rangle|
\le \|\mathcal{O}|\Psi(t)\rangle\|_2 \|\Psi(t)\|_2
\le \|\mathcal{O}\|_2 \|\Psi(t)\|_2^2.
\]
\end{proof}

Combining with the previous proposition yields:
\[
|\langle \Psi(t)|\mathcal{O}|\Psi(t)\rangle|
\le n\|\mathcal{O}\|_2\,\|A\|_2^2.
\]

This provides explicit operator-norm bounds for prime-labelled
tensor-network observables.

\nocite{*}
\bibliographystyle{unsrt}
\bibliography{references}

\end{document}